% Information Security & Acceptable Use Policy, officer variant.
%
% THE MOST IMPORTANT OF THE THREE OFFICER POLICIES, because an officer holds strictly more
% access than any intern will and previously had agreed to none of this.
%
% FOUR SECTIONS ARE SUBSTANTIVELY DIFFERENT. A word-swap would have produced a policy that
% is wrong about the one person it most needs to be right about:
%
%   5   Access          an intern may access only what an assignment requires. An officer
%                       legitimately holds broad and administrative access. Least privilege
%                       still applies, but as a discipline they impose on themselves rather
%                       than a boundary somebody else drew for them.
%   10  Security testing an intern may not authorise it. An officer can, and the rule that
%                       matters is that authorisation is recorded and scoped.
%   12  Incidents       org.yaml routes incidents to the CTO. An incident involving the CTO
%                       cannot be reported to the CTO.
%   15  Violations      an office is vacated or removed by the directors, not terminated.
%
% WORDING IS A DRAFT FOR REVIEW BY QUALIFIED COUNSEL.
\documentclass{paridyum}
\begin{document}

\paridyumletterhead
\paridyumdateref{\dkDocDate}{\dkDocRef}
\paridyumtitle{Information Security \& Acceptable Use Policy (Officer)}

\textit{This policy applies to officers of Paridyum, who are given access to its accounts,
source code, documents, systems, devices, infrastructure, or other digital resources.}

\section*{1. Purpose}
The purpose of this policy is to protect Paridyum's systems, source code, confidential
information, credentials, and other digital assets. An officer holds more access than
anyone else here, which is the reason this copy exists rather than the intern one.

\section*{2. Accounts and Passwords}
The Officer must:
\begin{itemize}[leftmargin=*, noitemsep]
    \item keep passwords and credentials private;
    \item use strong and unique passwords;
    \item enable multi-factor authentication where required;
    \item never share Company accounts with another person;
    \item never use another person's account without authorization; and
    \item immediately report lost, stolen, or compromised credentials.
\end{itemize}

\section*{3. Source Code and Repositories}
Paridyum's private source code must remain within authorized repositories and systems.
The Officer must not:
\begin{itemize}[leftmargin=*, noitemsep]
    \item publish private repositories;
    \item upload Company code to personal public repositories;
    \item share source code with unauthorized persons;
    \item copy Company code for unrelated personal projects; or
    \item intentionally remove or bypass repository security controls.
\end{itemize}

\section*{4. API Keys and Secrets}
Passwords, API keys, access tokens, private keys, database credentials, and other
secrets must never be:
\begin{itemize}[leftmargin=*, noitemsep]
    \item committed directly into source code;
    \item uploaded to public repositories;
    \item shared through unauthorized services;
    \item posted in public communication channels; or
    \item shared with another person without authorization.
\end{itemize}
If a secret is accidentally exposed, the Officer must report it immediately.

\section*{5. Access to Systems}
An officer's access is broad by necessity, and may be administrative. \textbf{The
discipline of least privilege therefore has to be self-imposed: being able to reach
something is not a reason to.} An officer accesses systems, repositories, databases,
files and information where the office requires it, and not otherwise --- and in
particular does not read another person's records because the permission happens to
exist.

The Officer must not:
\begin{itemize}[leftmargin=*, noitemsep]
    \item bypass or disable a security control rather than changing it openly;
    \item grant themselves access outside the process by which access is granted;
    \item use another person's account, or their credentials, in place of their own;
    \item view unrelated confidential information; or
    \item perform unauthorized security testing.
\end{itemize}

\section*{6. Personal Devices}
If a personal laptop or device is used for office work, the Officer should:
\begin{itemize}[leftmargin=*, noitemsep]
    \item use a password, PIN, or biometric device lock;
    \item keep the operating system reasonably updated;
    \item take reasonable precautions against malware;
    \item prevent unauthorized persons from accessing Company information; and
    \item securely remove Paridyum information when requested at the end of the
          office.
\end{itemize}

\section*{7. AI Tools}
The Officer must not provide confidential Paridyum information to external AI tools
unless their use has been approved. Do not submit the following to unauthorized AI
services:
\begin{itemize}[leftmargin=*, noitemsep]
    \item private source code;
    \item passwords or credentials;
    \item API keys;
    \item customer or user information;
    \item confidential documents;
    \item unreleased product information;
    \item internal business information; or
    \item sensitive technical architecture.
\end{itemize}
AI-generated code or content must be reviewed before being added to Paridyum projects.

\section*{8. Third-Party Services}
The Officer must not upload Paridyum information or source code to unauthorized:
\begin{itemize}[leftmargin=*, noitemsep]
    \item cloud-storage services;
    \item code-sharing websites;
    \item file-sharing services;
    \item online development environments;
    \item collaboration platforms; or
    \item other third-party services.
\end{itemize}
If a new service is required for the office, approval should be obtained before
using it with Paridyum information.

\section*{9. Software Installation}
Software used for Paridyum work should come from legitimate and reasonably trusted
sources. The Officer must not knowingly install or use:
\begin{itemize}[leftmargin=*, noitemsep]
    \item pirated software;
    \item malware;
    \item credential-stealing tools;
    \item unauthorized remote-access software; or
    \item software intended to bypass Paridyum's security controls.
\end{itemize}

\section*{10. Security Testing}
An officer may authorise security testing --- penetration testing, vulnerability
scanning, exploitation or load testing --- against Paridyum systems, and may carry it
out. Authorisation must be \textbf{recorded and scoped before the testing starts}, and
the testing must stay inside that scope. An officer does not authorise their own testing
of a system they alone administer without telling another officer it is happening.

\section*{11. Confidential Information}
Confidential information must be handled according to the \textbf{Paridyum
Non-Disclosure Agreement}. The Officer must take reasonable precautions to prevent
confidential information from being accidentally viewed, copied, published, or shared.

\section*{12. Security Incidents}
The Officer must immediately report suspected security incidents, including:
\begin{itemize}[leftmargin=*, noitemsep]
    \item leaked passwords or API keys;
    \item accidental public repositories;
    \item suspicious account activity;
    \item unauthorized access;
    \item malware;
    \item lost or stolen devices containing Paridyum information;
    \item accidental disclosure of confidential information; or
    \item discovered security vulnerabilities.
\end{itemize}

\vspace{0.3em}
\begin{infobox}
\textbf{Report security issues to}\\[0.4em]
\begin{tabularx}{\linewidth}{@{}>{\bfseries}p{2.2cm}X@{}}
Contact: & \dkOrgSecurityContact{} --- \dkOrgSecurityContactRole \\
Email:   & \dkOrgSecurityEmail
\end{tabularx}
\end{infobox}

\vspace{0.3em}
\textbf{Where the incident involves the security contact named above, it is reported to
another officer instead.} A report routed to the person it concerns is not a report, and
this is the one route in the policy that can point at its own subject.

\vspace{0.3em}
The Officer should report accidental incidents promptly rather than attempting to hide
or independently resolve a serious security issue.

\section*{13. Company Equipment}
Any equipment provided by Paridyum remains the property of Paridyum unless otherwise
agreed in writing. Company equipment must:
\begin{itemize}[leftmargin=*, noitemsep]
    \item be used responsibly;
    \item be reasonably protected from loss or damage;
    \item not be given to another person without permission; and
    \item be returned when requested or when the office ends.
\end{itemize}

\section*{14. Vacating the Office}
When the Officer ceases to hold office, they must:
\begin{itemize}[leftmargin=*, noitemsep]
    \item stop accessing Paridyum systems;
    \item return Company equipment;
    \item remove Company credentials from personal devices;
    \item return or delete confidential files when requested; and
    \item cooperate with account and access removal.
\end{itemize}
The Officer must not retain Paridyum's private source code, credentials, confidential
documents, or other protected information after the office.

\section*{15. Violations}
Serious or repeated violations of this policy may result in removal of system access, and
in the Officer being asked to vacate the office or removed from it by the directors. An
office is not terminated in the way an engagement is. Obligations under the
Non-Disclosure Agreement continue regardless, and where appropriate Paridyum may also
take other action available under applicable law.

\section*{16. Acknowledgement}
By signing below, the Officer confirms that they have read and understood this
Information Security \& Acceptable Use Policy and agree to follow it during their
office.

\vspace{1.5em}
\paridyumsignblock

\vspace{1.5em}
\paridyumcountersign{Officer}{\dkPIdentityLegalName}

\end{document}
